\section{From a dense set to every numerator}\label{sec:elementary}

The main construction produces short expansions for most numerators over a
carefully chosen denominator.  We first state exactly what is needed from
that construction, and then show why it suffices for every numerator.

\begin{proposition}[A dense set of short expansions]\label{prop:density}
There are absolute constants \(D_M>1\), \(L>0\), and \(S_0>1\) with the
following property.  Put \(D_X=D_M+2\) and \(D_C=4D_X\).
For every real \(S\ge S_0\) and every integer \(C\) satisfying
\[
 e^{D_CS}\le C\le e^{2D_CS},
\]
there are an integer \(M\) and a set
\(G\subseteq\{1,\ldots,\lfloor X\rfloor\}\), where \(X=e^{D_XS}\), such that
\[
 e^S<M\le e^{D_MS},
 \qquad
 \bigl|\{1,\ldots,\lfloor X\rfloor\}\setminus G\bigr|\le \frac X8.
\]
For each \(u\in G\), the fraction \(u/(MC)\) is a sum of at most
\(L\log S\) unit fractions with integer denominators at least \(2\).
Repetitions are allowed.
\end{proposition}

Section~\ref{sec:descent} proves Proposition~\ref{prop:density}.
Its constants and threshold are independent of \(C\); the integer \(M\)
and the set \(G\) may depend on \(C\).
The proof of Theorem~\ref{thm:main} from this proposition is elementary.
We record first how to remove repetitions and how to prepare a denominator
in the required range.

\begin{lemma}[Removing repetitions]\label{lem:distinct}
Let \(k\ge1\) be an integer.  If \(0<x<1\) is a sum of \(k\) positive unit fractions with integer
denominators, then it is a sum of \(k\) such fractions with distinct
denominators, all at least \(2\).
\end{lemma}

\begin{proof}
We use Takenouchi's argument \cite[Sections~II--III, pp.~79--80]{Takenouchi1921}.
For fixed \(x>0\) and \(k\), there are only finitely many nondecreasing
lists of \(k\) denominators with reciprocal sum \(x\).
Indeed, the first denominator is at most \(k/x\).
Once it is chosen, induction applies to the remaining positive sum and
the remaining \(k-1\) terms; when no terms remain the sum must be zero.

If a denominator occurs twice, apply one of the identities
\[
 \frac2{2p}=\frac1{p+1}+\frac1{p(p+1)},
 \qquad
 \frac2{2p+1}=\frac1{p+1}+\frac1{(p+1)(2p+1)}.
\]
Because the total is less than \(1\), a denominator \(1\), or a repeated
denominator \(2\), is impossible.  Thus \(p\ge2\) in the first identity
and \(p\ge1\) in the second.  Each replacement preserves the number of
terms and increases the sum of their denominators: the increases are
\((p-1)^2\) and \(2p^2\), respectively.  Finiteness of the set of lists
therefore forces termination, at which point all denominators are
distinct.  Positivity and the unchanged total exclude a denominator
\(1\) throughout.
\end{proof}

\begin{lemma}[Preparing the denominator]\label{lem:greedy}
Let \(1\le a<b\) be integers and let \(T>b\) be real.
After at most
\[
 1+\left\lceil\log_2\left(\frac{\log T}{\log 2}\right)\right\rceil
\]
greedy steps, either \(a/b\) has been expressed as a sum of unit fractions,
or its positive remainder can be written as \(A/C\), with integers
\[
 1\le A\le a,\qquad T\le C<T^2.
\]
All unit denominators produced are at least \(2\).
The fractions used during this procedure need not be reduced.
\end{lemma}

\begin{proof}
For a positive remainder \(A/C<1\), set
\[
 z=\left\lceil\frac CA\right\rceil,\qquad
 A'=Az-C,\qquad C'=Cz.
\]
Subtracting \(1/z\) leaves \(A'/C'\).
The ceiling inequality gives \(0\le A'<A\), while integrality of \(A\)
gives \(z\le C\) and hence \(C'\le C^2\).
Thus the numerator never exceeds \(a\).
Stop at zero or at the first positive remainder whose denominator is at
least \(T\).  At such a first crossing, the preceding denominator was
less than \(T\), so the new one is less than \(T^2\).

To bound the number of steps, write \(x=A/C\).  Since
\(\lceil1/x\rceil\le 1+1/x\),
\[
 0\le x-\frac1{\lceil1/x\rceil}
 \le \frac{x^2}{1+x}.
\]
The first positive remainder is less than \(1/2\), and subsequent
remainders decrease at least by squaring.  After \(j\ge1\) positive
steps their values therefore satisfy
\[
 x_j\le 2^{-2^{j-1}}.
\]
If the procedure has not yet stopped, its positive integer numerator
gives \(x_j\ge 1/C_j>1/T\).  This is impossible when
\(2^{j-1}\log2\ge\log T\), proving the stated bound.
Every current positive value is less than \(1\), so \(z\ge2\).
\end{proof}

\begin{proof}[Proof of Theorem~\ref{thm:main}, assuming
Proposition~\ref{prop:density}]
For the upper bound, fix \(1\le a<b\), put \(S=\log b\), and assume that
\(S\) is sufficiently large.  Apply Lemma~\ref{lem:greedy} with
\(T=e^{D_CS}\).  This uses \(O(\log S)\) terms.  If the procedure
terminates, Lemma~\ref{lem:distinct} gives the required upper bound.
Otherwise its
remainder \(A/C\) satisfies
\[
 1\le A<b,\qquad e^{D_CS}\le C<e^{2D_CS}.
\]
Choose \(M,G,X\) from Proposition~\ref{prop:density} for this \(S,C\).
Since \(D_X=D_M+2\),
\[
 AM<e^{(D_M+1)S}<X.
\]
Let \(g=\lfloor X/(AM)\rfloor\) and \(n=gAM\).  These are positive
integers, and the inequality \(\lfloor t\rfloor\ge t/2\) for \(t\ge1\)
gives
\[
 \frac X2\le n\le X.
\]
Write \(H=\{1,\ldots,\lfloor X\rfloor\}\setminus G\).
Among the \(n-1\) positive splits \(n=u+(n-u)\), at most \(2|H|\le X/4\)
have an entry outside \(G\).  Since \(n-1\ge X/2-1>X/4\) for large \(S\),
some split has both entries in \(G\).
Their expansions together express
\[
 \frac{n}{MC}=\frac{gA}{C}
\]
using at most \(2L\log S\) unit fractions.  Multiplying every denominator
by the integer \(g\) gives an expansion of \(A/C\) of the same length.
Adjoin the greedy prefix and apply Lemma~\ref{lem:distinct} to the total
\(a/b<1\).  The resulting distinct expansion has \(O(\log S)\) terms,
uniformly in \(a\).  No reduction of \(a/b\), or restriction on the final
denominator sizes, has been used.

For the lower bound, consider an expansion of \((b-1)/b\) with \(k\)
terms and append \(1/b\).  Sort the resulting \(s=k+1\) denominators as
\(d_1\le\cdots\le d_s\); repetitions are harmless here.
Put \(L_0=1\) and \(L_j=d_1\cdots d_j\).
For each \(1\le j\le s\), the amount remaining after the first \(j-1\)
terms is positive and has denominator dividing \(L_{j-1}\).
Consequently
\[
 \frac1{L_{j-1}}
 \le 1-\sum_{i<j}\frac1{d_i}
 =\sum_{i=j}^s\frac1{d_i}
 \le\frac{s}{d_j}.
\]
It follows that \(d_j\le sL_{j-1}\) and \(L_j\le sL_{j-1}^2\).
Induction gives \(L_j\le s^{2^j-1}\) and \(d_j\le s^{2^{j-1}}\).
As \(b\) is one of these denominators,
\[
 b\le d_s\le s^{2^{s-1}}=(k+1)^{2^k}.
\]
Hence
\(\log\log b\le k\log2+\log\log(k+1)\le(1+\log2)k\).
The numerator \(b-1\) therefore supplies the required lower bound for
\(N(b)\).
\end{proof}
